% problem_id: S001-proposition-nagata-universally-japanese
% source_id: S001 (Stacks Project)
% source_file: algebra.tex
% statement_locator: algebra.tex lines 46066--46074
% proof_locator: algebra.tex lines 46076--46222
% env: proposition
% status: text-extracted-verbatim (difficulty judgement pending, written by hand)
%
% The statement and proof below are the author's own text, extracted verbatim.
% Nothing is paraphrased, shortened, or supplemented.

\begin{proposition}[Nagata]
\label{proposition-nagata-universally-japanese}
Let $R$ be a ring. The following are equivalent:
\begin{enumerate}
\item $R$ is a Nagata ring,
\item any finite type $R$-algebra is Nagata, and
\item $R$ is universally Japanese and Noetherian.
\end{enumerate}
\end{proposition}

\begin{proof}
It is clear that a Noetherian universally Japanese ring is universally
Nagata (i.e., condition (2) holds). Let $R$ be a Nagata ring.
We will show that any finitely generated $R$-algebra $S$ is Nagata.
This will prove the proposition.

\medskip\noindent
Step 1. There exists a sequence of ring maps
$R = R_0 \to R_1 \to R_2 \to \ldots \to R_n = S$ such that
each $R_i \to R_{i + 1}$ is generated by a single element.
Hence by induction it suffices to prove $S$ is Nagata if
$S \cong R[x]/I$.

\medskip\noindent
Step 2. Let $\mathfrak q \subset S$ be a prime of $S$, and let
$\mathfrak p \subset R$ be the corresponding prime of $R$.
We have to show that $S/\mathfrak q$ is N-2. Hence we have
reduced to proving the following:
(*) Given a Nagata domain $R$ and a monogenic extension $R \subset S$
of domains then $S$ is N-2.

\medskip\noindent
Step 3: Let $R$ be a Nagata domain and
let $R \subset S$ be a monogenic extension of domains.
Suppose the induced extension of fraction fields of $R$ and $S$
is purely transcendental. In this case $S = R[x]$. By
Lemma \ref{lemma-polynomial-ring-N-2} we see that $S$ is N-2.
Hence we have reduced to proving the following:
(**) Given a Nagata domain $R$ and a
monogenic extension $R \subset S$ of domains
inducing a finite extension of fraction fields
then $S$ is N-2.

\medskip\noindent
Step 4. Let $R$ be a normal Nagata domain and
let $R \subset S$ be a monogenic extension of domains
inducing a finite extension of fraction fields $L/K$.
Choose an element $x \in S$
which generates $S$ as an $R$-algebra. Let $M/L$
be a finite extension of fields.
Let $R'$ be the integral closure of $R$ in $M$.
Then the integral closure $S'$ of $S$ in $M$ is equal to the integral
closure of $R'[x]$ in $M$.
Also the fraction field of $R'$ is $M$ and $R \subset R'$
is finite (by the Nagata property of $R$).
This implies that $R'$ is a Nagata ring
(Lemma \ref{lemma-quasi-finite-over-nagata}).
To show that $S'$ is finite over $S$ is the same as showing that
$S'$ is finite over $R'[x]$. Replace $R$ by $R'$ and $S$ by $R'[x]$
to reduce to the following statement:
(***) Given a normal Nagata domain $R$ with fraction field $K$,
and $x \in K$, the ring $S \subset K$ generated by $R$ and $x$
is N-1.

\medskip\noindent
Step 5. Let $R$ be a normal Nagata domain with fraction field $K$.
Let $x = b/a \in K$. We have to show that the ring $S \subset K$
generated by $R$ and $x$ is N-1. Note that $S_a \cong R_a$ is normal.
Hence by Lemma \ref{lemma-characterize-N-1} it suffices to show that
$S_{\mathfrak m}$ is N-1 for every maximal ideal $\mathfrak m$ of $S$.

\medskip\noindent
With assumptions as in the preceding paragraph, pick such a maximal
ideal and set $\mathfrak n = R \cap \mathfrak m$. The residue field
extension $\kappa(\mathfrak m)/\kappa(\mathfrak n)$ is finite
(Theorem \ref{theorem-nullstellensatz}) and generated by the image of $x$.
After replacing $R$ by $R_a$ for some $a \in R$, $a \not \in \mathfrak n$
and $S$ by $S_a$, may assume there exists a monic polynomial
$f(X) = X^d + \sum_{i = 1, \ldots, d} a_iX^{d - i}$ in $R[x]$ with
$f(x) \in \mathfrak m$ (details omitted).
Let $K''/K$ be a finite extension of fields such that the polynomial
$f(X)$ splits completely in $K''[X]$.
Let $R'$ be the integral closure of $R$ in $K''$.
Let $S' \subset K''$ be the subring generated by $R'$ and $x$.
As $R$ is Nagata we see $R'$ is finite over $R$ and Nagata
(Lemma \ref{lemma-quasi-finite-over-nagata}).
Moreover, $S'$ is finite over $S$. If for every maximal ideal
$\mathfrak m'$ of $S'$ the local ring $S'_{\mathfrak m'}$ is
N-1, then $S'_{\mathfrak m}$ is N-1 by
Lemma \ref{lemma-characterize-N-1}, which in turn
implies that $S_{\mathfrak m}$ is N-1 by
Lemma \ref{lemma-finite-extension-N-2}.
After replacing $R$ by $R'$ and $S$ by $S'$, and $\mathfrak m$ by
any of the maximal ideals $\mathfrak m'$ lying over $\mathfrak m$
we reach the situation where the polynomial $f$ above split completely:
$f(X) = \prod_{i = 1, \ldots, d} (X - a_i)$ with $a_i \in R$.
Since $f(x) \in \mathfrak m$ we see that $x - a_i \in \mathfrak m$
for some $i$. Finally, after replacing $x$ by $x - a_i$ we may assume
that $x \in \mathfrak m$.

\medskip\noindent
To recapitulate: $R$ is a normal Nagata domain with fraction field $K$,
$x \in K$ and $S$ is the subring of $K$ generated by $x$ and $R$,
finally $\mathfrak m \subset S$ is a maximal ideal with $x \in \mathfrak m$.
We have to show $S_{\mathfrak m}$ is N-1.

\medskip\noindent
We will show that Lemma \ref{lemma-criterion-analytically-unramified}
applies to the local ring
$S_{\mathfrak m}$ and the element $x$. This will imply that
$S_{\mathfrak m}$ is analytically unramified, whereupon we
see that it is N-1 by Lemma \ref{lemma-analytically-unramified-easy}.

\medskip\noindent
We have to check properties (1), (2), (3)(a) and (3)(b).
Property (1) is trivial.
Let $I = \Ker(R[X] \to S)$ where $X \mapsto x$.
We claim that $I$ is generated by all linear forms $aX - b$ such that
$ax = b$ in $K$. Clearly all these linear forms are in $I$.
If $g = a_d X^d + \ldots a_1 X + a_0 \in I$, then we see that
$a_dx$ is integral over $R$ (Lemma \ref{lemma-make-integral-trivial})
and hence $b := a_dx \in R$
as $R$ is normal. Then $g - (a_dX - b)X^{d - 1} \in I$ and we win by
induction on the degree. As a consequence we see that
$$
S/xS = R[X]/(X, I) = R/J
$$
where
$$
J = \{b \in R \mid ax = b \text{ for some }a \in R\} = xR \cap R
$$
By Lemma \ref{lemma-normal-domain-intersection-localizations-height-1}
we see that $S/xS = R/J$ has no embedded primes as an $R$-module, hence as
an $R/J$-module, hence as an $S/xS$-module, hence as an $S$-module.
This proves property (2).
Take such an associated prime $\mathfrak q \subset S$ with the
property $\mathfrak q \subset \mathfrak m$ (so that it is an
associated prime of $S_{\mathfrak m}/xS_{\mathfrak m}$ -- it does not
matter for the arguments).
Then $\mathfrak q$ is minimal over $xS$ and hence has height $1$.
By the sequence of equalities above we see that
$\mathfrak p = R \cap \mathfrak q$ is an associated
prime of $R/J$, and so has height $1$
(see Lemma \ref{lemma-normal-domain-intersection-localizations-height-1}).
Thus $R_{\mathfrak p}$ is a discrete valuation ring and therefore
$R_{\mathfrak p} \subset S_{\mathfrak q}$ is an equality. This shows
that $S_{\mathfrak q}$ is regular. This proves property (3)(a).
Finally, $(S/\mathfrak q)_{\mathfrak m}$ is a localization
of $S/\mathfrak q$, which is a quotient of $S/xS = R/J$.
Hence $(S/\mathfrak q)_{\mathfrak m}$ is a localization of
a quotient of the Nagata ring $R$, hence
Nagata (Lemmas \ref{lemma-quasi-finite-over-nagata}
and \ref{lemma-nagata-localize})
and hence analytically unramified
(Lemma \ref{lemma-local-nagata-domain-analytically-unramified}).
This shows (3)(b) holds and we are done.
\end{proof}
